% ============================================================
% Iterazione/caso studio 1
% ============================================================

\chapter{Iterazione 1: SOLIDWORKS + STAR-CCM+ + HEEDS}
\label{cap:iter1-solidworks}

\section{Descrizione del caso}

Il primo caso studio usa SOLIDWORKS come CAD parametrico, STAR-CCM+ come ambiente di preparazione geometria, mesh e CFD, e HEEDS come orchestratore di ottimizzazione. La pipeline e':

\begin{center}
\texttt{SOLIDWORKS -> Parasolid/STEP -> STAR-CCM+ mesh/CFD -> HEEDS -> SOLIDWORKS}
\end{center}

Questo workflow e' credibile per un team che progetta gia' in SOLIDWORKS e vuole sfruttare STAR-CCM+ per il calcolo CFD industriale. La parte critica e' mantenere stabile il modello CAD durante molte varianti automatiche.

\section{Workflow operativo}

\begin{enumerate}
    \item Creare una configurazione CAD dedicata alla simulazione, semplificata rispetto al modello produttivo.
    \item Esporre in SOLIDWORKS solo parametri robusti: raggi principali, sezioni, lunghezze, posizioni e angoli.
    \item Esportare preferibilmente Parasolid \texttt{.x\_t} o STEP, evitando STL come formato CAD principale.
    \item In STAR-CCM+ importare la geometria, riparare superfici, definire regioni e boundary names.
    \item Usare automated meshing con controlli locali e prism layers.
    \item Estrarre reports: drag, lift, momenti, residual monitors, eventuali proxy di power/thrust.
    \item HEEDS aggiorna i parametri SOLIDWORKS, rilancia la catena e salva ogni variante.
\end{enumerate}

\section{Punti di forza}

\begin{itemize}
    \item Curva di ingresso piu' bassa se il team conosce SOLIDWORKS.
    \item Ottimo per parti meccaniche, assiemi, bracket, carenature e disegni produttivi.
    \item STAR-CCM+ compensa molte fragilita' downstream grazie a surface repair, wrapper e meshing automatico.
    \item HEEDS permette di esplorare molte varianti senza riscrivere manualmente ogni caso.
\end{itemize}

\section{Rischi}

\begin{itemize}
    \item Feature tree fragile: fillet o split line possono fallire al variare dei parametri.
    \item Naming superfici instabile: se le patch cambiano nome, le boundary conditions possono rompersi.
    \item SOLIDWORKS e' meno naturale di NX per superfici aerospace complesse, compositi e digital thread Simcenter.
    \item Workflow proprietario costoso: non puo' essere la sola base replicabile della repo.
\end{itemize}

\section{Valutazione}

\begin{table}[H]
\centering
\caption{Valutazione del caso SOLIDWORKS + STAR-CCM+ + HEEDS}
\begin{tabular}{P{5cm} c c}
\toprule
\textbf{Criterio} & \textbf{Voto} & \textbf{Nota} \\
\midrule
Analisi supportate & 5 & STAR-CCM+ copre RANS/URANS, moving mesh e multiphysics \\
Accuratezza & 4.5 & Dipende da mesh, modello fisico e validazione \\
Integrabilita' & 4 & Buona, ma meno nativa di NX nel portfolio Siemens \\
Costo licenze & 1 & Stack commerciale completo \\
Scalabilita' & 5 & STAR-CCM+/HEEDS supportano workflow distribuiti \\
Compatibilita' OS & 4 & SOLIDWORKS vincola a Windows per il CAD \\
Supporto formati & 4 & Usare Parasolid/STEP come scambio principale \\
\midrule
\textbf{Totale} & \textbf{3.90/5} & \textbf{Valido, ma non ideale come benchmark proprietario principale} \\
\bottomrule
\end{tabular}
\end{table}

\section{Decisione}

SOLIDWORKS + STAR-CCM+ + HEEDS e' consigliabile se il team e' gia' SOLIDWORKS-first. Per il campo droni/eVTOL e' una soluzione operativa valida, ma richiede CAD governance rigorosa: template, naming standard, design table controllate, export pulito e regression test su rigenerazione geometrica.
